\documentclass{article}
\usepackage[UTF8]{ctex}
\usepackage{amsmath,amssymb,amsfonts}
\usepackage{enumitem}
\usepackage{geometry}
\geometry{a4paper, margin=1in}
% === 错题本排版优化（不参与题目标识） ===
\setlength{\parindent}{0pt}
\setlength{\parskip}{0.35em plus 0.10em minus 0.08em}
\setlength{\emergencystretch}{3em}
\setlist{itemsep=0.25em, topsep=0.40em, parsep=0pt}
\allowdisplaybreaks[2]
% === 错题本排版优化结束 ===
\begin{document}
\section*{北京邮电大学2024年数学分析真题}
\section*{14．北京邮电大学2024年数学分析试卷}
一．求极限

\[
\lim _{x \rightarrow 0} \frac{1-\cos x \sqrt{\cos 2 x} \sqrt[3]{\cos 3 x}}{x^{2}}
\]

\begin{itemize}
\end{itemize}

\subsection{题目2}
原题边界“二．”：计算
\[
\int_{-\frac{\pi}{2}}^{\frac{\pi}{2}} \frac{1+\sin 3 x}{1+\cos ^{2} x} \mathrm{~d} x
\]

三．求幂级数 \(\sum_{n=1}^{\infty} \frac{(-1)^{(n-1)}}{2 n-1} x^{2 n}\) 的收敛域及和函数．

四．求二元函数 \(f(x, y)=x^{2}\left(2+y^{2}\right)+y \ln y\) 的极值．

五．若 \(a_{n}>0, \lim _{n \rightarrow \infty} \frac{a_{n}}{a_{n+1}}=l>1\) ，证明： \(\lim _{n \rightarrow \infty} a_{n}=0\) ．

六．讨论 \(\ln ^{2} x-2 \ln x+2 x+k=0\) 的实根个数，其中 \(k\) 是常数．

七．已知 \(x \mathrm{e}^{-\frac{x}{2}}-\ln (1+x)\) 与 \(\alpha x^{n}\) 是等价无穷小量，求 \(\alpha\) 和 \(n\) ．

八．判别 \(\int_{0}^{+\infty} \frac{\ln x}{\sqrt[3]{x^{4}-x}} \mathrm{~d} x\) 的敛散性。

九．求圆锥体 \(z \geq \sqrt{x^{2}+y^{2}}\) 与 \(x^{2}+y^{2}+(z-1)^{2} \leq 1\) 所围立体的体积．

十．计算

\[
I=\int_{L} 3 x y \mathrm{~d} x+\left(x^{3}+x-2 y\right) \mathrm{d} y
\]

其中 \(L\) 是第一象限从 \((0,0)\) 沿圆周 \(x^{2}+y^{2}-2 x=0\) 到（2，0），再从（2，0）沿圆周 \(x^{2}+y^{2}=4\) 到（0，2）的曲线．

十一．平面 \(2 x+y+2 z=2\) 与三个坐标平面围成 \(\Omega, \sum\) 为 \(\Omega\) 的整个表面，取外侧，求

\[
I=\iint_{\Sigma}\left(x^{2}+1\right) \mathrm{d} y \mathrm{~d} z-2 y \mathrm{~d} z \mathrm{~d} x+3 z \mathrm{~d} x \mathrm{~d} y
\]

十二．讨论 \(\alpha>0\) 时，\(f(x)=x^{\alpha}\) 在 \((0,+\infty)\) 上的一致连续性．

十三。证明 \(\sum_{n=1}^{\infty} \frac{(-1)^{n}}{x^{2}+n}\) 在 \((-\infty,+\infty)\) 上一致收敛。

十四．\(f(x)\) 在 \((0,+\infty)\) 上可微， \(0 \leq f(x) \leq \frac{x}{1+x^{2}}\) ，证明存在 \(\xi>0\) ，使 \(f^{\prime}(\xi)=\frac{1-\xi^{2}}{\left(1+\xi^{2}\right)^{2}}\) ．

十五。设 \(f(x)\) 可导且 \(f(0)=1,0<f^{\prime}(x)<\frac{1}{2}\) ，数列 \(\left\{x_{n}\right\}\) 满足 \(x_{n+1}=f\left(x_{n}\right)(n=1,2, \cdots)\) ．证明：

\begin{itemize}

  \item[1．] 级数 \(\sum_{n=1}^{\infty}\left(x_{n+1}-x_{n}\right)\) 绝对收敛。

  \item[2．] \(\lim _{n \rightarrow \infty} x_{n}\) 存在且 \(0<\lim _{n \rightarrow \infty} x_{n}<2\) ．
\end{itemize}
\end{document}
